% =====================================================================
% PREÂMBULO PADRÃO — ROTEIROS PRÁTICOS (visual moderno, desde o Roteiro 02)
% UC Iniciação Científica e Extensionista — IFSC Garopaba
% Documentação: .agents/base-de-conhecimento/padrao-visual-roteiros-pdf.md
%
% Como usar em um novo roteiro:
%   1. Copie este preâmbulo para o .tex do roteiro (do \documentclass até
%      antes de \begin{document}) — ou use \input se preferir;
%   2. Ajuste \numeroRoteiro e \tituloCurto logo abaixo;
%   3. Copie figures/garopaba_horizontal_marca2015_PNG.png para a pasta do roteiro;
%   4. Compile com: tectonic arquivo.tex
% =====================================================================
\documentclass[11pt,a4paper]{article}

\newcommand{\numeroRoteiro}{0X}                 % <-- número do roteiro
\newcommand{\tituloCurto}{Título curto do roteiro} % <-- aparece no cabeçalho

% =====================================================
% PACOTES
% =====================================================
\usepackage[utf8]{inputenc}
\usepackage[T1]{fontenc}
\usepackage[brazil]{babel}
\shorthandoff{"}
\usepackage[default]{sourcesanspro}          % texto: sem serifa, leitura confortável
\usepackage[scale=0.86]{sourcecodepro}       % código: monoespaçada moderna
\usepackage{fontawesome5}                    % ícones
\usepackage{geometry}
\usepackage{graphicx}
\usepackage{fancyhdr}
\usepackage{titlesec}
\usepackage[table]{xcolor}
\usepackage{hyperref}
\usepackage{enumitem}
\usepackage[most]{tcolorbox}
\usepackage{float}
\usepackage{tikz}
\usetikzlibrary{arrows.meta, positioning, shapes.geometric, calc, babel}
\usepackage{tabularx}
\usepackage{booktabs}
\usepackage{listings}
\usepackage{lastpage}
\usepackage{qrcode}
\usepackage{multicol}

\geometry{top=2.6cm, bottom=2.6cm, left=2cm, right=2cm, headsep=0.8cm}
\linespread{1.18}
\setlength{\parskip}{0.55em}
\setlength{\parindent}{0pt}

% =====================================================
% PALETA (cores institucionais IFSC, em tons suaves)
% =====================================================
\definecolor{azul}{RGB}{30,60,90}
\definecolor{azulMedio}{RGB}{37,99,235}
\definecolor{verde}{RGB}{22,101,52}
\definecolor{verdeVivo}{RGB}{22,163,74}
\definecolor{ambar}{RGB}{217,119,6}
\definecolor{vermelho}{RGB}{220,38,38}
\definecolor{roxo}{RGB}{124,58,237}
\definecolor{tinta}{RGB}{30,41,59}
\definecolor{suave}{RGB}{100,116,139}
\definecolor{borda}{RGB}{226,232,240}
\definecolor{fundo}{RGB}{248,250,252}
\definecolor{barraTitulo}{RGB}{241,245,249}
\definecolor{terminal}{RGB}{15,23,42}

\color{tinta}

% =====================================================
% LISTINGS
% =====================================================
\lstdefinestyle{base}{
    basicstyle=\ttfamily\small\color{tinta},
    keywordstyle=\color{azulMedio}\bfseries,
    commentstyle=\color{suave}\itshape,
    stringstyle=\color{verde},
    breaklines=true, columns=fullflexible, keepspaces=true,
    showstringspaces=false, upquote=true, tabsize=4,
    frame=none, numbers=none, aboveskip=0pt, belowskip=0pt,
    extendedchars=true,
    literate=
      {á}{{\'a}}1 {é}{{\'e}}1 {í}{{\'i}}1 {ó}{{\'o}}1 {ú}{{\'u}}1
      {Á}{{\'A}}1 {É}{{\'E}}1 {Í}{{\'I}}1 {Ó}{{\'O}}1 {Ú}{{\'U}}1
      {ã}{{\~a}}1 {õ}{{\~o}}1 {Ã}{{\~A}}1 {Õ}{{\~O}}1
      {â}{{\^a}}1 {ê}{{\^e}}1 {ô}{{\^o}}1 {Â}{{\^A}}1 {Ê}{{\^E}}1 {Ô}{{\^O}}1
      {à}{{\`a}}1 {ç}{{\c c}}1 {Ç}{{\c C}}1
      {–}{{--}}1 {±}{{$\pm$}}1 {≈}{{$\approx$}}1 {÷}{{$\div$}}1
}
\lstdefinestyle{py}{style=base, language=Python,
    morekeywords={True,False,None}}
\lstdefinestyle{tex}{style=base, language=TeX}
\lstdefinestyle{txt}{style=base}
\lstdefinestyle{saida}{style=base, language={}, basicstyle=\ttfamily\small\color{white!88!black}}
\lstset{style=base}

% =====================================================
% COMPONENTES VISUAIS
% =====================================================
% --- Janela de código (estilo editor) ---------------
\newcommand{\pontosJanela}{%
  {\color{vermelho!70}\tiny\faCircle}\,{\color{ambar!80}\tiny\faCircle}\,{\color{verdeVivo!70}\tiny\faCircle}\hspace{0.8em}}
\newtcblisting{codigo}[2][py]{
    enhanced, listing only, breakable,
    colback=fundo, colframe=borda, boxrule=0.6pt, arc=6pt,
    left=10pt, right=10pt, top=6pt, bottom=6pt,
    title={\pontosJanela #2}, coltitle=suave, colbacktitle=barraTitulo,
    fonttitle=\ttfamily\footnotesize, toptitle=2pt, bottomtitle=2pt,
    titlerule=0.6pt, titlerule style=borda,
    listing options={style=#1}
}
% --- Saída (terminal escuro) -------------------------
\newtcblisting{saida}[1][Saída]{
    enhanced, listing only, breakable,
    colback=terminal, colframe=terminal, boxrule=0pt, arc=6pt,
    left=10pt, right=10pt, top=4pt, bottom=6pt,
    title={\faTerminal\hspace{0.5em}#1}, coltitle=white!65!black,
    colbacktitle=terminal, fonttitle=\sffamily\footnotesize\bfseries,
    toptitle=3pt, bottomtitle=0pt, titlerule=0pt,
    listing options={style=saida}
}
% --- Exemplo certo/errado -----------------------------
\newtcblisting{exemplo}[2]{
    enhanced, listing only,
    colback=#1!4, colframe=#1!45, boxrule=0.6pt, arc=6pt,
    left=8pt, right=8pt, top=5pt, bottom=5pt,
    title={#2}, colbacktitle=#1!12, coltitle=#1!75!black,
    fonttitle=\sffamily\small\bfseries, toptitle=2pt, bottomtitle=2pt,
    listing options={style=py, basicstyle=\ttfamily\footnotesize\color{tinta}}
}
% --- Caixas de destaque com barra lateral -----------
\newtcolorbox{destaque}[3]{
    enhanced, breakable,
    colback=#1!6, frame hidden, arc=6pt,
    borderline west={3.5pt}{0pt}{#1},
    left=14pt, right=12pt, top=9pt, bottom=9pt,
    before upper={{\color{#1}\large #2}\hspace{0.55em}{\color{#1!75!black}\bfseries #3}\par\vspace{0.25em}}
}
\newenvironment{dica}[1]{\begin{destaque}{ambar}{\faLightbulb[regular]}{#1}}{\end{destaque}}
\newenvironment{atencao}[1]{\begin{destaque}{vermelho}{\faExclamationTriangle}{#1}}{\end{destaque}}
\newenvironment{conceito}[1]{\begin{destaque}{azulMedio}{\faBookOpen}{#1}}{\end{destaque}}
\newenvironment{sucesso}[1]{\begin{destaque}{verdeVivo}{\faCheckCircle[regular]}{#1}}{\end{destaque}}
\newenvironment{avancado}[1]{\begin{destaque}{roxo}{\faCode}{#1}}{\end{destaque}}

% --- Caixa com título "aba" (objetivos, atividade) ---
\newtcolorbox{painel}[3][]{
    enhanced, breakable,
    colback=white, colframe=borda, boxrule=0.6pt, arc=8pt,
    left=14pt, right=14pt, top=14pt, bottom=10pt,
    title={#2\hspace{0.5em}#3}, coltitle=white, colbacktitle=azul,
    fonttitle=\bfseries, attach boxed title to top left={xshift=12pt, yshift=-8pt},
    boxed title style={arc=4pt, boxrule=0pt}, #1
}
% --- Cartões pequenos ---------------------------------
\tcbset{cartao/.style={enhanced, colback=fundo, colframe=borda, boxrule=0.6pt, arc=8pt,
    left=10pt, right=10pt, top=8pt, bottom=8pt, fontupper=\small\raggedright}}

% --- Elementos em linha --------------------------------
\newcommand{\tecla}[1]{\tcbox[on line, enhanced, colback=white, colframe=suave!50, boxrule=0.5pt,
    arc=2.5pt, left=3pt, right=3pt, top=0.5pt, bottom=0.5pt, boxsep=0.5pt,
    borderline south={1.2pt}{0pt}{suave!40}]{\ttfamily\footnotesize #1}}
\newcommand{\botao}[1]{\tcbox[on line, enhanced, colback=azulMedio!8, colframe=azulMedio!35, boxrule=0.5pt,
    arc=3pt, left=3pt, right=3pt, top=0.5pt, bottom=0.5pt, boxsep=0.5pt]{\sffamily\footnotesize\color{azul}#1}}
\newcommand{\seta}{\,{\color{suave}\scriptsize\faChevronRight}\,}
\newcommand{\chip}[2]{\tcbox[on line, enhanced, colback=white, opacityback=0.14, colframe=white,
    opacityframe=0.45, boxrule=0.5pt, arc=8pt, left=5pt, right=5pt, top=1.5pt, bottom=1.5pt,
    boxsep=0pt, coltext=white]{\footnotesize #1\hspace{0.4em}#2}}
\newcommand{\circulo}[2][azul]{\tikz[baseline=(c.base)]\node[circle, fill=#1, text=white,
    inner sep=0pt, minimum size=1.35em, font=\footnotesize\bfseries](c){#2};}
\newcommand{\marca}[1]{\tikz[baseline=(c.base)]\node[circle, fill=vermelho, text=white,
    inner sep=0pt, minimum size=1.2em, font=\scriptsize\bfseries](c){#1};}
\newcommand{\caixaMarcar}{{\color{suave}\faSquare[regular]}}

% =====================================================
% HIPERLINKS
% =====================================================
\hypersetup{colorlinks=true, linkcolor=azulMedio, urlcolor=verde, citecolor=azul}

% =====================================================
% CABEÇALHO E RODAPÉ
% =====================================================
\pagestyle{fancy}
\fancyhf{}
\fancyhead[L]{%
    \IfFileExists{figures/garopaba_horizontal_marca2015_PNG.png}{%
        \raisebox{-0.25\height}{\includegraphics[height=0.95cm]{figures/garopaba_horizontal_marca2015_PNG.png}}%
    }{\textbf{\color{verde} IFSC — Câmpus Garopaba}}}
\fancyhead[R]{\footnotesize\color{suave}\textbf{\color{azul}Roteiro Prático \numeroRoteiro}\hspace{0.5em}\textbullet\hspace{0.5em}\tituloCurto}
\fancyfoot[L]{\footnotesize\color{suave}Iniciação Científica e Extensionista\hspace{0.4em}\textbullet\hspace{0.4em}CST Sistemas para a Internet}
\fancyfoot[R]{\tcbox[on line, colback=azul!8, colframe=azul!8, arc=6pt, boxrule=0pt, left=5pt, right=5pt, top=1pt, bottom=1pt, boxsep=0pt]{\footnotesize\color{azul}\bfseries\thepage\,/\,\pageref*{LastPage}}}
\renewcommand{\headrule}{{\color{borda}\hrule height 0.6pt}}
\renewcommand{\footrulewidth}{0pt}
\setlength{\headheight}{36pt}

% =====================================================
% TÍTULOS DE SEÇÃO
% =====================================================
\titleformat{\section}
  {\normalfont\Large\bfseries\color{azul}}
  {\tikz[baseline=(n.base)]\node[fill=azul, rounded corners=5pt, text=white, inner sep=5pt,
      minimum width=1.9em, font=\large\bfseries](n){\thesection};}
  {0.7em}{}
\titlespacing*{\section}{0pt}{1.6em}{0.7em}

\titleformat{\subsection}
  {\normalfont\large\bfseries\color{azul!85!black}}
  {{\color{verdeVivo}\thesubsection}}
  {0.6em}{}
\titlespacing*{\subsection}{0pt}{1.2em}{0.4em}

\setlist[itemize]{leftmargin=1.4em, itemsep=0.25em, topsep=0.3em, label={\color{verdeVivo}\small\faAngleRight}}
\setlist[enumerate]{leftmargin=2.2em, itemsep=0.45em, topsep=0.3em}

% Tabelas: sem linhas verticais, cabeçalho colorido, zebra suave
\newcommand{\cabecalho}{\rowcolor{azul}}
\newcommand{\cab}[1]{\leavevmode\color{white}\bfseries #1}
\arrayrulecolor{borda}
\newcolumntype{Y}{>{\raggedright\arraybackslash}X}
\newcolumntype{P}[1]{>{\raggedright\arraybackslash}p{#1}}

% \begin{document} ... (conteúdo do roteiro)
